% Input from reports/phase2/ with amsmath and
% cite loaded. New keys: sections/literature_review.bib (consolidate in #80);
% Phase 1 keys reused from reports/phase1/refs.bib.
% Every number in Table~\ref{tab:lit_dl} was read from the cited paper's own
% table (paper, table): Chen2021machine IV; Jin2022bilstm III; Zhang2022dual III;
% Gupta2026auditing 5; Galchar2026remaining III; Sharma2025uncertainty I.

\section{Deep Learning on C-MAPSS}
\label{sec:lit_dl}

Phase~1 reviewed feature-based regression on C-MAPSS and noted that
convolutional and recurrent networks report lower errors than those
baselines~\cite{SateeshBabu2016deep,Li2018remaining,Zheng2017long}. This
section looks at the four families of sequence models that Phase~2 implements,
the numbers they report on FD001, and why those numbers cannot simply be
ranked.

\subsection{Recurrent Networks}
An LSTM~\cite{Hochreiter1997long} reads a window of sensor readings one cycle
at a time and keeps a gated memory, so it can carry a slow wear trend across
the window. Zheng et al.~\cite{Zheng2017long} established the
window-plus-LSTM recipe on C-MAPSS. Chen et al.~\cite{Chen2021machine} kept
that recipe (window 30, ten runs averaged) and asked what an LSTM is missing.
In their ablation, the plain LSTM reached an FD001 RMSE of 15.43. Adding
attention over time steps lowered it to 15.27, adding handcrafted features
lowered it to 14.81, and doing both lowered it to 14.54. Handcrafted features
helped more than attention. Jin et al.~\cite{Jin2022bilstm} pushed the same
idea further, with a bidirectional LSTM on the raw sensors (12.58) and a second
stream on handcrafted feature flows, fused by addition (11.96).

\subsection{Convolutional and Temporal Convolutional Networks}
A 1D convolution slides learned filters along the time axis of a window, so it
picks up local patterns, such as a short rise in temperature, in parallel
rather than step by step~\cite{SateeshBabu2016deep,Li2018remaining}. A
temporal convolutional network (TCN)~\cite{Bai2018empirical} makes the
convolutions causal (each output sees only the past) and dilated (gaps between
filter taps grow layer by layer), so its receptive field grows exponentially
with depth. That lets a few layers cover a whole window. In a single-seed
comparison, Galchar et al.~\cite{Galchar2026remaining} report an FD001 RMSE of
16.97 for a 1D CNN and 14.93 for an LSTM under the same protocol. TCN variants
also appear in comparison tables of later papers~\cite{Zhang2022dual}, but we
could not check a TCN result on FD001 against an open primary source, so none
is quoted here.

\subsection{Transformers}
Self-attention~\cite{Vaswani2017attention} lets every cycle in the window look
at every other cycle directly, instead of passing information along a chain.
Mo et al.~\cite{Mo2021remaining} combine a Transformer encoder with a gated
convolutional unit for local features. Zhang et al.~\cite{Zhang2022dual}
(DAST) use two attention encoders, one across sensors and one across time
steps, and report an FD001 RMSE of 11.43 (window 40, ten runs). Transformers
have many parameters for a training set of only 100 engines, so these results
depend heavily on regularization and on the protocol details below.

\subsection{Hybrid Models}
The strongest recurrent results above are already hybrids. They pair a learned
sequence representation with handcrafted
features~\cite{Chen2021machine,Jin2022bilstm}, which is the idea behind our
gated hybrid. Sharma~\cite{Sharma2025uncertainty} stacks multi-scale
convolutions, a bidirectional LSTM and attention, and adds a probabilistic
output that predicts its own uncertainty. That model reports 17.25 on FD001 as
an average over runs and 16.22 for its best run.

\begin{table*}[t]
\centering
\caption{FD001 test RMSE reported by deep and hybrid models. Every value was
read from the paper's own results table. Cap: training RUL cap. Test capped:
whether test labels were also capped. ``n.s.'': not stated. $^\dagger$Preprint.}
\label{tab:lit_dl}
\footnotesize
\setlength{\tabcolsep}{2.5pt}
\begin{tabular}{llccccc}
\hline
Paper & Model & RMSE & Cap & Test capped & Window & Runs \\
\hline
Chen~\cite{Chen2021machine} & LSTM & 15.43 & n.s.$^a$ & n.s. & 30 & 10 \\
 & LSTM + features & 14.81 & & & & \\
 & + attention & 14.54$^b$ & & & & \\
Jin~\cite{Jin2022bilstm} & Bi-LSTM & 12.58 & 125 & n.s. & 30 & n.s. \\
 & two-stream & 11.96 & & & & \\
Zhang (DAST)~\cite{Zhang2022dual} & DAST & 11.43 & 125 & n.s. & 40$^c$ & 10 \\
Gupta~\cite{Gupta2026auditing}$^\dagger$ & LSTM & 15.97$\pm$0.45 & 125 & no & 30 & 5 seeds \\
Galchar~\cite{Galchar2026remaining}$^\dagger$ & LSTM & 14.93 & 130 & n.s. & 30 & 1 seed \\
 & 1D CNN & 16.97 & & & & \\
 & Ridge & 16.63 & & & & \\
Sharma~\cite{Sharma2025uncertainty}$^\dagger$ & CNN-BiLSTM-att. & 17.25$^d$ & 125 & n.s. & 30 & n.s. \\
\hline
Ours, Phase 1 & Random Forest & 16.21 & 125 & no & --- & 1 seed \\
\hline
\end{tabular}

\smallskip
\raggedright\scriptsize
$^a$Piecewise-linear target; the maximum is not given.
$^b$Table~IV; the same model is 14.53 in their Table~V.
$^c$4 of the 100 FD001 test engines are shorter than 40 cycles; padding is not described.
$^d$Mean over runs; best run 16.22.
\end{table*}

\subsection{Why the Numbers Disagree}
The spread in Table~\ref{tab:lit_dl} (11.4 to 17.3 for broadly similar models)
is not only about architecture. Six protocol choices move the result.
\begin{enumerate}
\item \emph{Cap and test labels.} The training cap changes what the model
learns. Capping the test labels as well changes the exam itself: 11 of the
100 FD001 test engines have a true RUL above 125 (the maximum is 145), and a
capped label erases most of their error. Several papers state the cap for
training but do not say how test labels were treated.
\item \emph{Window length.} Longer windows give more history, but FD001's
shortest test engine has only 31 cycles, so any window above 31 needs padding,
which is rarely described. Chen et al.~\cite{Chen2021machine} found FD001
RMSE was best at 30 and worse for both shorter and longer windows.
\item \emph{Seeds and reporting.} Gupta~\cite{Gupta2026auditing} measures a
seed standard deviation of 0.45 RMSE for an FD001 LSTM. Gaps between papers
smaller than about one cycle are therefore within run-to-run noise. Reporting
a single seed, or the best run instead of the mean~\cite{Sharma2025uncertainty},
makes a result look better than a repeat would.
\item \emph{Model selection on the test set.} C-MAPSS has no official
validation split, so when a paper tunes architecture or epochs on the test
engines, the test RMSE becomes optimistic. Papers rarely say which split
was used for tuning.
\item \emph{Copied baselines.} Comparison tables often copy baseline numbers
from earlier papers, each with its own protocol~\cite{Zhang2022dual,Jin2022bilstm}.
Even one paper's own tables can disagree in the last digit (14.54 vs.\ 14.53
in~\cite{Chen2021machine}).
\item \emph{Preprocessing.} Sensor selection and normalization matter little
on FD001, which has one operating condition, but they dominate on FD002 and
FD004~\cite{Gupta2026auditing}. In Gupta's grid, architecture explains 75\% of
the FD001 RMSE variation, while on FD002 and FD004 the largest shares go to
the RUL cap and normalization.
\end{enumerate}
These are the same comparability problems that Ramasso and
Saxena~\cite{Ramasso2014performance} described for the earlier literature.

\subsection{Where Our Protocol Sits}
We train on FD001 with a cap of 125, chosen from training engines only
(Section~\ref{sec:target}). The official test labels stay uncapped, as
provided. Windows are 30 cycles long, with first-cycle padding. All model and
hyperparameter choices are made with engine-grouped cross-validation on the
training engines, and the official test set is opened once for the frozen
models. This places us closest to Gupta's audit~\cite{Gupta2026auditing}: the
same cap, uncapped test labels, the same window and several seeds. His LSTM at
$15.97\pm0.45$ and our Phase~1 capped Random Forest at 16.21 are therefore the
fair reference points for our deep models. The lower values in
Table~\ref{tab:lit_dl} come from protocols that are partly unstated, and
should not be treated as a target that a correct implementation must reach.
